\documentclass{beamer}

\usepackage[utf8]{inputenc}
\input{../helper}

\title{Introdução à Computação \\ árvores binárias de busca}
\author{Alexandre Rademaker}
\date{}

\begin{document}

\begin{frame}
 \maketitle
\end{frame}


\begin{frame}[allowframebreaks,fragile]{árvores}

  Uma árvore é um tipo de dados que representa um conjunto de nós
  conectados em uma estrutura hierárquica.

  Cada nó de uma árvore pode ser conectado a um número qualquer de
  \emph{filhos} - dependendo do tipo de árvore - mas apenas à um
  \emph{pai}. Um único nó, chamado \emph{raiz}, não tem pai.

  Isto siginifica que árvores não podem conter \emph{ciclos} e que
  cada nó pode ser considerado raiz de uma sub-árvore.

  Com isso, \emph{recursão} é uma técnica útil para travesia da
  árvore.

  \framebreak

  \begin{columns}
    \begin{column}{.4\textwidth}
      \begin{itemize}
      \item raiz
      \item folhas
      \item sub-árvores
      \item grau
      \item nível
      \item altura
      \end{itemize}
    \end{column}
    \begin{column}{.6\textwidth}
      \begin{center}
        \includegraphics[width=.9\textwidth]{trees.png}
      \end{center}
    \end{column}
  \end{columns}
  
  \framebreak

  \emph{Árvores binárias} são um tipo de árvore onde cada nó pode
  conter até dois filhos.

  Quando a ordem dos filhos é relevante, chamamos de \emph{árvore
    ordenada}.

  Valores podem ser associados a cada nó na árvore ou, as vezes,
  apenas aos nós \emph{folhas} - nós que não tem filhos.

  \framebreak

  aplicações:

  \begin{itemize}
  \item sistema de arquivos dos sistemas operacionais
  \item DOM de páginas HTML e arquivos XML
  \item árvores de busca para armazenamento de dados de forma a
    permitir buscas eficientes.
  \item árvores sintáticas no processamento de linguagem natural.
  \end{itemize}

  \framebreak

  Um \emph{tipo abstrato de dados} pode ser representado de diferentes
  formas.

  Na memória, tipicamente nós são alocados dinamicamente com ponteiros
  para os filhos, pais, ou ambos, além dos dados associados ao nó.

  Se o número de filhos é variado, a cada nó podemos usar um array ou
  lista encadeada para armazenar ponteiros para os filhos.

  Também podemos representar a árvore com duas listas, uma para os nós
  e outra para as relações pai-filho (lista de adjacencias), ou um
  array onde os relacionamentos (pai-filho) são determinados pelas
  posições no array.

  \framebreak

  \begin{lstlisting}[style=code,language=lisp]
    '(2 (7 (2) (10) (6 (5) (11))) 
        (5 (9 (4))))
  \end{lstlisting}

  Se nossa implementação de lista encadeada permitisse que elementos
  da lista fossem de qualquer tipo, incluíndo listas\ldots

  também poderíamos implementar uma árvore como uma lista de listas,
  onde o primeiro elemento de cada sublista é a raiz da sub-árvore.

  Para saber \href{https://en.wikipedia.org/wiki/Tree_(data_structure)}{mais}!

\end{frame}


\begin{frame}[allowframebreaks, fragile]{árvores binárias de busca}

  Uma árvore binária de busca, também chamada de árvore binária
  ordenada, é uma árvore binária onde o valor de cada nó interno é
  maior que todos os valores na subárvore esquerda do respectivo nó e
  menor que as de sua subárvore direita.

  A complexidade de tempo das operações na árvore de busca binária é
  diretamente proporcional à altura da árvore.

  Buscas binárias podem ser facilmente implementadas em árvores
  binárias de busca.

  \framebreak

  \includegraphics[width=.6\textwidth]{btree.png}
  
  \framebreak
  
  Em C, os nós precisam ser agora declarados com

  \begin{lstlisting}[style=normalc,caption=src/tree.c]
    typedef struct node {
      int key;
      struct node *left, *right;
    } t_node;
  \end{lstlisting}

  \framebreak

  Para inserir um novo elemento, uma função auxiliar que cria um novo
  nó:

  \begin{lstlisting}[style=normalc,caption=src/tree.c]
    t_node* new_node(int item) {
      t_node* t = malloc(sizeof(struct node));
      assert(t != NULL);
      t->key = item;
      t->left = t->right = NULL;
      return t;
    }
  \end{lstlisting}

  Notação \texttt{a = b = expressão}!

  \framebreak

  Agora a inserção, seja o nó corrente com valor V, devemos comparar
  com o valor X sendo inserido. Se X menor, movemos para a árvore à
  esquerda. Se X maior, para árvore da direita.

  \begin{lstlisting}[style=normalc,caption=src/tree.c]
    t_node* insert(t_node* node, int key) {
      if (node == NULL)
        return new_node(key);
      
      if (key < node->key)
        node->left = insert(node->left, key);
      else if (key > node->key)
        node->right = insert(node->right, key);
      
      return node;
    }
  \end{lstlisting}

  \framebreak

  \includegraphics[width=.8\textwidth]{btree-insert.png}

  \framebreak

  Nós criados de forma dinâmica precisam ser liberados.

  \begin{lstlisting}[style=normalc, caption=src/tree.c]
    void free_tree(node *root) {
      if (root == NULL)
        return;

      free_tree(root->left);
      free_tree(root->right);
      free(root);
    }
  \end{lstlisting}

\end{frame}

\begin{frame}[allowframebreaks,fragile]{busca binária}

  Usando esta estrutura, podemos então realizar buscas\vspace{.5em}

  \begin{lstlisting}[style=normalc,caption=src/tree.c]
    bool binary_search(t_node *root, int n) {
      if (root == NULL)
        return false;
      else if (n < root->key)
        return binary_search(root->left, n);
      else if (n > root->key)
        return binary_search(root->right, n);
      return (n == root->key);
    }
  \end{lstlisting}

  \framebreak

  \begin{columns}
    \begin{column}{.5\textwidth}
      Mas temos que 'balancear' a árvore, para não terminar com uma
      estruta ineficiente.\vspace{1em}

      Em uma árvore balanceada, busca e inserção são $O(\log n)$, mas
      em uma árvore não balanceada, pode ser $O(n)$.
    \end{column}
    \begin{column}{.5\textwidth}
      \begin{center}
        \includegraphics[width=.8\textwidth]{imbalanced_tree.png}
      \end{center}
    \end{column}
  \end{columns}
    
\end{frame}

\end{document}

%%% Local Variables:
%%% mode: latex
%%% TeX-master: t
%%% End:
